% !TeX program = pdflatex
\documentclass[a4paper]{article}

% --- Gói ngôn ngữ và font ---
\usepackage[T5,T1]{fontenc}    % Vietnamese support with T1 as the main encoding
\usepackage[utf8]{inputenc}    % UTF-8 source text
\usepackage{lmodern}           % Vector fonts for pdfLaTeX
\usepackage[english]{babel}    % Hỗ trợ tiếng Anh (tùy chọn)
\usepackage{amsmath}
\usepackage{amssymb}
\usepackage{latexsym}
\usepackage{indentfirst} % Ép buộc thụt đầu dòng đoạn đầu tiên

% --- Gói hình ảnh và màu sắc ---
\usepackage{graphicx}
\usepackage{xcolor} % Dùng xcolor thay cho color để mạnh mẽ hơn
\usepackage{tikz}
\usetikzlibrary{arrows,backgrounds}
\usepackage[framemethod=tikz]{mdframed}

% --- Gói định dạng trang và bảng ---
\usepackage{geometry}
\usepackage{a4wide}
\usepackage{fancyhdr}
\usepackage{booktabs}
\usepackage{array}
\usepackage{hhline}
\usepackage{multicol}
\usepackage{multirow}
\usepackage{makecell}
\usepackage{lastpage}
\usepackage{placeins}
\usepackage{setspace}
\usepackage[skip=0pt]{caption}
\usepackage{float} % Để cố định hình ảnh với [H]

% --- Gói code và thuật toán ---
\usepackage{listings}
\usepackage[most]{tcolorbox}

% --- Các gói khác ---
\usepackage{blindtext}
\usepackage{comment}
\usepackage{enumerate}
\usepackage{enumitem}
\usepackage{hyperref}

% --- Cấu hình link ---
\hypersetup{
    urlcolor=blue,
    linkcolor=black,
    citecolor=black,
    colorlinks=true
}

% --- Định nghĩa màu sắc ---
\setlength{\columnseprule}{1pt}
\def\columnseprulecolor{\color{black}}
\definecolor{codegreen}{rgb}{0,0.6,0}
\definecolor{codegray}{rgb}{0.5,0.5,0.5}
\definecolor{codepurple}{rgb}{0.58,0,0.82}
\definecolor{backcolour}{rgb}{0.95,0.95,0.92}
\definecolor{ReportNavy}{HTML}{102A43}
\definecolor{ReportCyan}{HTML}{32B5D2}
\definecolor{CodeBackground}{HTML}{F6F8FA}
\definecolor{CodeFrame}{HTML}{CBD5E1}
\definecolor{CodeComment}{HTML}{2E7D32}
\definecolor{CodeKeyword}{HTML}{005CC5}
\definecolor{CodeString}{HTML}{B31D28}
\definecolor{CodeNumber}{HTML}{64748B}

% Optional logo asset: place your official hcmut.png alongside this file.
\newcommand{\reportlogo}[1]{%
    \IfFileExists{hcmut.png}{\includegraphics[width=#1]{hcmut.png}}{%
        \ifdim#1<1cm
          \begin{tikzpicture}[baseline=(mark.base)]
          \node[draw=ReportNavy,fill=ReportNavy,text=white,rounded corners=1pt,
                minimum width=#1,minimum height=#1,
                inner sep=0pt,font=\tiny\bfseries] (mark) {H};
          \end{tikzpicture}%
        \else
          \begin{tikzpicture}[baseline=(mark.base)]
          \node[draw=ReportNavy,fill=ReportNavy,text=white,rounded corners=2pt,
                minimum width=#1,minimum height=18mm,
                inner sep=4pt,font=\Large\bfseries] (mark) {HCMUT};
          \end{tikzpicture}%
        \fi
    }%
}

% --- Cấu hình Minted (C++) ---
% Minted is not needed; tcolorbox/listings supplies the template code boxes.
%\usemintedstyle{vs}
% --- Cấu hình Header/Footer ---
\everymath{\color{blue}}
\setlength{\headheight}{40pt}
\pagestyle{fancy}
\fancyhead{} % clear all header fields
\fancyhead[L]{
    \begin{tabular}{rl}
        \begin{picture}(25,15)(0,0)
            \put(0,0){\reportlogo{8mm}}
        \end{picture}&
        \begin{tabular}{l}
            \textbf{\ttfamily University of Technology}\\
            \textbf{\ttfamily Faculty of Computer Science and Engineering}
        \end{tabular} 	
    \end{tabular}
}
\fancyhead[R]{} % Bên phải header để trống

\fancyfoot{} % clear all footer fields
\fancyfoot[L]{\scriptsize \ttfamily \textcolor{blue}{Microprocessors--Microcontrollers (CO3010) -- Lab 02 Timer}}
\fancyfoot[R]{\scriptsize \ttfamily Page {\thepage}/\pageref{LastPage}}
\renewcommand{\headrulewidth}{0.3pt}
\renewcommand{\footrulewidth}{0.3pt}

% --- Cấu hình Mục lục ---
\setcounter{secnumdepth}{4}
\setcounter{tocdepth}{3}
\makeatletter
\newcounter {subsubsubsection}[subsubsection]
\renewcommand\thesubsubsubsection{\thesubsubsection .\@alph\c@subsubsubsection}
\newcommand\subsubsubsection{\@startsection{subsubsubsection}{4}{\z@}%
                                     {-3.25ex\@plus -1ex \@minus -.2ex}%
                                     {1.5ex \@plus .2ex}%
                                     {\normalfont\normalsize\bfseries}}
\newcommand*\l@subsubsubsection{\@dottedtocline{3}{10.0em}{4.1em}}
\newcommand*{\subsubsubsectionmark}[1]{}
\makeatother

\setlength{\parindent}{1cm}
\setlength{\parskip}{0pt}

\begin{document}

% --- TRANG BÌA ---
\hypersetup{pageanchor=false}
\begin{titlepage}
\begin{center}
    \Large Vietnam National University, Ho Chi Minh City \\
    \Large University of Technology \\
    \Large Faculty of Computer Science and Engineering
\end{center}

\vspace{1cm}

\begin{figure}[H]
    \centering
    \reportlogo{4cm}
\end{figure}

\vspace{1cm}

\begin{center}
    \begin{tabular}{c}
    \multicolumn{1}{l}{\textbf{{\Large Microprocessors--Microcontrollers (Lab) (CO3010)}}}\\
    ~~\\
    \hline
    \\
    \multicolumn{1}{l}{\textbf{{\Large Laboratory Report}}}\\
    \\
    \textbf{{\Large LAB\_02: TIMER INTERRUPT}}\\
    \textbf{{\Large AND LED SCANNING}}\\
    \\
    \hline
    \end{tabular}
\end{center}

\vspace{2cm}
\begin{table}[H]
    \centering
    \begin{tabular}{rrl}
        & Instructor(s): & Phan Văn Sỹ, CSE-HCMUT \\
        & Class: & CC03\\
        & Student: & Vo Le Thanh Trieu -- 2453297\\
    \end{tabular}
\end{table}

\vspace{1cm}
\begin{center}
    {\footnotesize Ho Chi Minh City, October 2026}
\end{center}
\end{titlepage}
\hypersetup{pageanchor=true}

\thispagestyle{fancy}
\newpage
\tableofcontents
\newpage

% ------------------------------------------------------------
% Astonishing STM32 C code block
% ------------------------------------------------------------
\newtcblisting[auto counter,number within=section]{stmcode}[1]{
    enhanced,
    breakable,
    listing only,
    colback=CodeBackground,
    colframe=CodeFrame,
    colbacktitle=ReportNavy,
    coltitle=white,
    fonttitle=\bfseries,
    title={Program~\thetcbcounter: #1},
    attach boxed title to top left={xshift=4mm,yshift=-2mm},
    boxed title style={
        colback=ReportNavy,
        colframe=ReportNavy,
        arc=1.5mm,
        boxrule=0pt
    },
    arc=2mm,
    boxrule=0.7pt,
    left=7mm,
    right=2mm,
    top=4mm,
    bottom=2mm,
    drop shadow southeast,
    before skip=1.2em,
    after skip=1.2em,
    listing options={
        language=C,
        basicstyle=\ttfamily\small,
        keywordstyle=\color{CodeKeyword}\bfseries,
        commentstyle=\color{CodeComment}\itshape,
        stringstyle=\color{CodeString},
        numbers=left,
        numberstyle=\tiny\color{CodeNumber},
        numbersep=10pt,
        stepnumber=1,
        showstringspaces=false,
        breaklines=true,
        breakatwhitespace=false,
        keepspaces=true,
        tabsize=4,
        columns=fullflexible,
        escapeinside={(*@}{@*)}
    }
}

% ------------------------------------------------------------
% Reusable figure placeholder
% The real image is loaded automatically when the file exists.
% ------------------------------------------------------------
\newcommand{\schematicplaceholder}[2]{%
    \IfFileExists{#1}{%
        \includegraphics[width=#2\textwidth]{#1}%
    }{%
        \fcolorbox{ReportCyan}{CodeBackground}{%
            \parbox[c][5.5cm][c]{0.82\textwidth}{%
                \centering
                \textcolor{ReportNavy}{\Large\bfseries Schematic Placeholder}\\[0.8em]
                Add the image file:\\[0.4em]
                \texttt{\detokenize{#1}}
            }%
        }%
    }%
}

% ============================================================
\section{Timer Interrupt and LED Scanning}

This report follows Exercises 1--10 of \emph{Lab 02: Timer Interrupt and LED
Scanning}.

STM32CubeMX configures the STM32F103C6; application code is edited
and built in Visual Studio Code with CMake. Proteus provides the circuit
simulation. The code listings document successive versions of the same
application: later exercises extend the earlier ones. Screenshots of the
student's own Proteus simulations can be inserted into the named figure files.


% ============================================================
% Replace the existing Exercise 1 subsection with this section.
% The original figure path, caption, and label are preserved.
% ============================================================
\subsection{Exercise 1: Two seven-segment displays}

\subsubsection{Objective and GPIO configuration}

Exercise 1 uses timer interrupts to control two common-anode
seven-segment displays. The first display shows the number 1,
and the second shows the number 2. Only one display is enabled
at a time, and the selected display changes every 500~ms.

The segment cathodes of both displays share PB0--PB6.
Separate PNP transistors control their common anodes through
PA6 and PA7. This arrangement allows the two displays to
share the same segment signals.

\begin{table}[H]
    \centering
    \caption{GPIO configuration for Exercise 1}
    \label{tab:ex01-gpio}
    \begin{tabular}{lll}
        \toprule
        \textbf{GPIO pins} &
        \textbf{Signals} &
        \textbf{Function} \\
        \midrule
        PB0--PB6 &
        SEG\_A--SEG\_G &
        Shared segment outputs, active low \\
        PA6 &
        EN0 &
        First display enable, active low \\
        PA7 &
        EN1 &
        Second display enable, active low \\
        \bottomrule
    \end{tabular}
\end{table}

PB0, PB1, PB2, PB3, PB4, PB5, and PB6 correspond to
segments a, b, c, d, e, f, and g, respectively.
These pins are configured as push-pull outputs with no
internal pull-up or pull-down resistors.

For the selected common-anode display, writing
\texttt{GPIO\_PIN\_RESET} to a segment output turns that
segment on. Writing \texttt{GPIO\_PIN\_SET} turns it off.
The PNP enable signals are also active low: PA6 or PA7
must be low to enable its corresponding display.

\subsubsection{Schematic}

\begin{figure}[H]
    \centering
    \schematicplaceholder{figures/pic02.png}{0.86}
    \caption{Proteus schematic for Exercise 1: two common-anode digits and PNP enables.}
    \label{fig:ex01}
\end{figure}

\subsubsection{Step-by-step timer calculation}

The timer configuration follows the introductory timer
setup in the lab manual. TIM2 generates an interrupt every
10~ms, and the application counts these interrupts to obtain
the required 500-ms display-selection interval.

\paragraph{Step 1: Determine the TIM2 input clock.}

The source code selects the internal HSI oscillator,
disables the PLL, and sets the AHB and APB1 clock dividers
to 1. Therefore,
\[
    f_{\mathrm{SYSCLK}}
    = f_{\mathrm{HCLK}}
    = f_{\mathrm{PCLK1}}
    = 8\,000\,000~\mathrm{Hz}.
\]

Because the APB1 divider is 1, the TIM2 input clock equals
the APB1 peripheral clock:
\[
    f_{\mathrm{TIM2}}
    = 8\,000\,000~\mathrm{Hz}
    = 8~\mathrm{MHz}.
\]

\paragraph{Step 2: Calculate the counter clock after the prescaler.}

The source code sets
\[
    PSC = 7999.
\]

The prescaler divides the timer input clock by
$PSC+1$, so the counter frequency is
\[
    f_{\mathrm{CNT}}
    = \frac{f_{\mathrm{TIM2}}}{PSC+1}
    = \frac{8\,000\,000}{7999+1}
    = \frac{8\,000\,000}{8000}
    = 1000~\mathrm{Hz}.
\]

Thus, the counter advances once every
\[
    T_{\mathrm{CNT}}
    = \frac{1}{f_{\mathrm{CNT}}}
    = \frac{1}{1000}~\mathrm{s}
    = 1~\mathrm{ms}.
\]

\paragraph{Step 3: Calculate the timer update interval.}

TIM2 uses up-counting mode, and the source code sets
\[
    ARR = 9.
\]

The counter counts from 0 through 9. An update occurs
when it overflows and returns to 0, giving
\[
    N_{\mathrm{counts}} = ARR+1 = 10
\]
counter ticks per update.

Therefore,
\[
    T_{\mathrm{interrupt}}
    = (ARR+1)T_{\mathrm{CNT}}
    = 10 \times 1~\mathrm{ms}
    = 10~\mathrm{ms}.
\]

Equivalently,
\[
    T_{\mathrm{interrupt}}
    = \frac{(PSC+1)(ARR+1)}{f_{\mathrm{TIM2}}}
    = \frac{8000 \times 10}{8\,000\,000}
    = 0.01~\mathrm{s}.
\]

\paragraph{Step 4: Calculate the interrupt frequency.}

The timer interrupt frequency is
\[
    f_{\mathrm{interrupt}}
    = \frac{1}{T_{\mathrm{interrupt}}}
    = \frac{1}{0.01}
    = 100~\mathrm{Hz}.
\]

With TIM2 interrupts enabled, the callback is therefore
invoked 100 times per second.

\paragraph{Step 5: Calculate the number of interrupts per display.}

Each display must remain selected for
\[
    T_S = 500~\mathrm{ms}.
\]

The required number of timer interrupts is
\[
    N_{\mathrm{interrupts}}
    = \frac{T_S}{T_{\mathrm{interrupt}}}
    = \frac{500~\mathrm{ms}}{10~\mathrm{ms}}
    = 50.
\]

This explains the condition \texttt{ticks >= 50}
in the callback. The hardware timer still interrupts
every 10~ms, while the display changes every
\[
    50 \times 10~\mathrm{ms}
    = 500~\mathrm{ms}.
\]

\subsubsection{Segment-display function}

The actual \texttt{display7SEG()} implementation in
Exercise 1 handles only the numbers 1 and 2.

Before writing a new digit, the function sets PB0--PB6
high to turn off all segments. It then drives the
required segment outputs low:
\begin{itemize}
    \item Number 1 uses segments b and c:
          PB1 and PB2.
    \item Number 2 uses segments a, b, d, e, and g:
          PB0, PB1, PB3, PB4, and PB6.
\end{itemize}

The prototype is placed in \texttt{USER CODE BEGIN PFP},
and the function definition is placed in
\texttt{USER CODE BEGIN 0}.

\begin{stmcode}{Exercise 1 display7SEG function}
/* USER CODE BEGIN PFP */
void display7SEG(int num);
/* USER CODE END PFP */

/* USER CODE BEGIN 0 */
void display7SEG(int num){
    /* Turn off all seven segments first */
    HAL_GPIO_WritePin(
        GPIOB,
        GPIO_PIN_0 | GPIO_PIN_1 | GPIO_PIN_2 | GPIO_PIN_3 |
        GPIO_PIN_4 | GPIO_PIN_5 | GPIO_PIN_6,
        GPIO_PIN_SET
    );

    if (num == 1)
    {
        /* Number 1 uses segments b and c */
        HAL_GPIO_WritePin(
            GPIOB,
            GPIO_PIN_1 | GPIO_PIN_2,
            GPIO_PIN_RESET
        );
    }
    else if (num == 2)
    {
        /* Number 2 uses segments a, b, d, e, and g */
        HAL_GPIO_WritePin(
            GPIOB,
            GPIO_PIN_0 | GPIO_PIN_1 | GPIO_PIN_3 |
            GPIO_PIN_4 | GPIO_PIN_6,
            GPIO_PIN_RESET
        );
    }
  }
/* USER CODE END 0 */
\end{stmcode}

\subsubsection{Initialization in main}

After GPIO and TIM2 initialization, both displays are
disabled by setting PA6 and PA7 high. The segment pattern
for number 1 is then prepared, and PA6 is driven low to
enable the first display.

Finally, \texttt{HAL\_TIM\_Base\_Start\_IT(\&htim2)}
starts TIM2 in interrupt mode. The TIM2 interrupt must
also be enabled in the CubeMX NVIC settings.

The following excerpt matches the initialization
in the supplied source code.

\begin{stmcode}{Exercise 1 initialization inside main}
MX_GPIO_Init();
MX_TIM2_Init();

/* USER CODE BEGIN 2 */

/* Disable both displays */
HAL_GPIO_WritePin(
    GPIOA,
    GPIO_PIN_6 | GPIO_PIN_7,
    GPIO_PIN_SET
);

/* Prepare number 1, then enable the first display */
display7SEG(1);
HAL_GPIO_WritePin(GPIOA, GPIO_PIN_6, GPIO_PIN_RESET);

/* Start TIM2 with interrupts */
HAL_TIM_Base_Start_IT(&htim2);

/* USER CODE END 2 */
\end{stmcode}

The \texttt{while (1)} loop contains no application
processing in this exercise. After initialization,
display switching is handled by the timer callback,
without calling \texttt{HAL\_Delay()}.

\subsubsection{Timer interrupt callback}

The callback checks \texttt{htim->Instance} so that
the display logic executes only for TIM2.

Two static variables retain their values between
callback invocations:
\begin{itemize}
    \item \texttt{ticks} counts the 10-ms timer events.
    \item \texttt{active\_display} records which display
          is currently enabled. Value 0 represents
          the first display, and value 1 represents
          the second.
\end{itemize}

The initial value of \texttt{active\_display} is 0,
matching the first display enabled during startup.
Every callback increments \texttt{ticks}. Once it
reaches 50, the counter is reset and the other
display is selected.

\begin{stmcode}{Exercise 1 TIM2 callback}
/* USER CODE BEGIN 4 */
void HAL_TIM_PeriodElapsedCallback(TIM_HandleTypeDef *htim){
    if (htim->Instance == TIM2)
    {
        static uint8_t ticks = 0;
        static uint8_t active_display = 0;

        ticks++;

        if (ticks >= 50)
        {
            ticks = 0;

            /* Turn off both displays before changing segments */
            HAL_GPIO_WritePin(
                GPIOA,
                GPIO_PIN_6 | GPIO_PIN_7,
                GPIO_PIN_SET
            );

            if (active_display == 0)
            {
                /* Change from the first display to the second */
                display7SEG(2);
                HAL_GPIO_WritePin(
                    GPIOA, GPIO_PIN_7, GPIO_PIN_RESET
                );
                active_display = 1;
            }
            else
            {
                /* Change from the second display to the first */
                display7SEG(1);
                HAL_GPIO_WritePin(
                    GPIOA, GPIO_PIN_6, GPIO_PIN_RESET
                );
                active_display = 0;
            }
        }
    }
  }
/* USER CODE END 4 */
\end{stmcode}

Each switching operation follows three steps:
disable both displays, write the new segment pattern,
and enable the selected display. Disabling the displays
before changing the shared segment signals prevents
the previous display from briefly showing the new pattern.

The callback executes once for each timer event and
then returns. It is not an infinite loop. Between
interrupts, the GPIO outputs retain their levels,
so the selected display remains enabled.

\subsubsection{Expected simulation behavior}

Immediately after initialization, the first display
shows 1 and the second is disabled. After 50 timer
events, approximately 500~ms later, the second
display shows 2 and the first is disabled.
The sequence then repeats.

\begin{table}[H]
    \centering
    \caption{Nominal display sequence for Exercise 1}
    \label{tab:ex01-sequence}
    \begin{tabular}{llll}
        \toprule
        \textbf{Time interval} &
        \textbf{PA6 / EN0} &
        \textbf{PA7 / EN1} &
        \textbf{Visible digit} \\
        \midrule
        0--0.5~s   & LOW  & HIGH & First display: 1 \\
        0.5--1.0~s & HIGH & LOW  & Second display: 2 \\
        1.0--1.5~s & LOW  & HIGH & First display: 1 \\
        1.5--2.0~s & HIGH & LOW  & Second display: 2 \\
        \bottomrule
    \end{tabular}
\end{table}

Because the required selection interval is 500~ms,
the alternating displays are clearly visible in
Proteus. This slow scan demonstrates the multiplexing
process; it does not make both digits appear
continuously illuminated.

\subsubsection{Short question: Frequency of the scanning process}

The introduction defines the scanning frequency
using the time required to visit every display.
For $N$ displays, each selected for an interval $T_S$,
\[
    T_{\mathrm{scan}} = N T_S,
    \qquad
    f_{\mathrm{scan}} = \frac{1}{N T_S}.
\]

Exercise 1 has two displays, with
$T_S = 500~\mathrm{ms} = 0.5~\mathrm{s}$. Therefore,
\[
    T_{\mathrm{scan}}
    = 2 \times 0.5~\mathrm{s}
    = 1~\mathrm{s},
\]
and
\[
    \boxed{
        f_{\mathrm{scan}}
        = \frac{1}{1~\mathrm{s}}
        = 1~\mathrm{Hz}
    }.
\]

A different display is selected every 0.5~s,
giving a selection rate of
\[
    f_{\mathrm{selection}}
    = \frac{1}{0.5~\mathrm{s}}
    = 2~\mathrm{Hz}.
\]

Thus, the TIM2 interrupt frequency is 100~Hz,
the display-selection rate is 2~Hz, and the
complete scanning frequency is \textbf{1~Hz}.

% ============================================================
% ============================================================
% Replace the existing Exercise 2 subsection with this section.
% The original figure path, caption, and label are preserved.
% ============================================================
\subsection{Exercise 2: Four digits and the center dot}

\subsubsection{Objective and GPIO configuration}

Exercise 2 extends Exercise 1 by adding two more
common-anode seven-segment displays and two center
LEDs controlled by the DOT signal on PA4.

The four displays show the digits 1, 2, 3, and 0,
representing 12:30. Only one display is enabled
at a time, and the selected display changes every
500~ms. Both DOT LEDs change state together
every second.

The segment outputs remain connected to PB0--PB6.
PA6--PA9 control the four PNP digit-enable circuits.

\begin{table}[H]
    \centering
    \caption{GPIO configuration for Exercise 2}
    \label{tab:ex02-gpio}
    \begin{tabular}{lll}
        \toprule
        \textbf{GPIO pins} &
        \textbf{Signals} &
        \textbf{Function} \\
        \midrule
        PB0--PB6 & SEG\_A--SEG\_G & Shared segment outputs \\
        PA6 & EN0 & First display: digit 1 \\
        PA7 & EN1 & Second display: digit 2 \\
        PA8 & EN2 & Third display: digit 3 \\
        PA9 & EN3 & Fourth display: digit 0 \\
        PA4 & DOT & Two center LEDs \\
        \bottomrule
    \end{tabular}
\end{table}

The segment outputs, digit enables, and DOT signal
are active low. A selected digit is enabled by
driving its enable pin low, while the other
enable pins remain high.

\subsubsection{Schematic}

\begin{figure}[H]
    \centering
    \schematicplaceholder{figures/pic01.png}{0.86}
    \caption{Proteus schematic for Exercise 2: four digits, DOT, and D1.}
    \label{fig:ex02}
\end{figure}

\subsubsection{Step-by-step timing calculation}

TIM2 retains the introductory configuration:
an 8-MHz input clock, prescaler 7999, and
auto-reload value 9.

\paragraph{Step 1: Calculate the timer interrupt interval.}

The counter frequency after the prescaler is
\[
    f_{\mathrm{CNT}}
    = \frac{8\,000\,000}{7999+1}
    = 1000~\mathrm{Hz}.
\]

Thus, each counter tick lasts
\[
    T_{\mathrm{CNT}}
    = \frac{1}{1000}~\mathrm{s}
    = 1~\mathrm{ms}.
\]

Since the counter counts from 0 through 9,
each update requires ten counter ticks:
\[
    T_{\mathrm{interrupt}}
    = (9+1)\times 1~\mathrm{ms}
    = 10~\mathrm{ms}.
\]

The corresponding interrupt frequency is
\[
    f_{\mathrm{interrupt}}
    = \frac{1}{0.01~\mathrm{s}}
    = 100~\mathrm{Hz}.
\]

\paragraph{Step 2: Calculate the digit-selection threshold.}

Each display must remain selected for 500~ms.
Therefore, the number of interrupts required
before selecting the next display is
\[
    N_{\mathrm{scan}}
    = \frac{500~\mathrm{ms}}{10~\mathrm{ms}}
    = 50.
\]

This gives the callback condition
\texttt{scan\_ticks >= 50}.

\paragraph{Step 3: Calculate the DOT-toggle threshold.}

DOT must change state every second:
\[
    T_{\mathrm{DOT\ toggle}} = 1000~\mathrm{ms}.
\]

The required number of interrupts is
\[
    N_{\mathrm{DOT}}
    = \frac{1000~\mathrm{ms}}{10~\mathrm{ms}}
    = 100.
\]

This gives the callback condition
\texttt{dot\_ticks >= 100}.

The two counters advance during every TIM2 callback,
but each counter resets only when its own threshold
is reached.

\subsubsection{Segment-display function}

The supplied Exercise 2 source implements
\texttt{display7SEG()} for digits 0, 1, 2, and 3.
These are the four values required to display 12:30.

Before selecting a new pattern, the function sets
all segment outputs high to turn them off.
It then drives the required segment outputs low.

\begin{table}[H]
    \centering
    \caption{Segment patterns used in Exercise 2}
    \label{tab:ex02-patterns}
    \begin{tabular}{cl}
        \toprule
        \textbf{Digit} & \textbf{Segments turned on} \\
        \midrule
        0 & a, b, c, d, e, f \\
        1 & b, c \\
        2 & a, b, d, e, g \\
        3 & a, b, c, d, g \\
        \bottomrule
    \end{tabular}
\end{table}

\begin{stmcode}{Exercise 2 display7SEG function}
/* USER CODE BEGIN PFP */
void display7SEG(int num);
/* USER CODE END PFP */

/* USER CODE BEGIN 0 */
void display7SEG(int num)
{
    /* First turn all seven segments off */
    HAL_GPIO_WritePin(
        GPIOB,
        GPIO_PIN_0 | GPIO_PIN_1 | GPIO_PIN_2 | GPIO_PIN_3 |
        GPIO_PIN_4 | GPIO_PIN_5 | GPIO_PIN_6,
        GPIO_PIN_SET
    );

    if (num == 0)
    {
        /* 0: a, b, c, d, e, f */
        HAL_GPIO_WritePin(
            GPIOB,
            GPIO_PIN_0 | GPIO_PIN_1 | GPIO_PIN_2 |
            GPIO_PIN_3 | GPIO_PIN_4 | GPIO_PIN_5,
            GPIO_PIN_RESET
        );
    }
    else if (num == 1)
    {
        /* 1: b, c */
        HAL_GPIO_WritePin(
            GPIOB,
            GPIO_PIN_1 | GPIO_PIN_2,
            GPIO_PIN_RESET
        );
    }
    else if (num == 2)
    {
        /* 2: a, b, d, e, g */
        HAL_GPIO_WritePin(
            GPIOB,
            GPIO_PIN_0 | GPIO_PIN_1 | GPIO_PIN_3 |
            GPIO_PIN_4 | GPIO_PIN_6,
            GPIO_PIN_RESET
        );
    }
    else if (num == 3)
    {
        /* 3: a, b, c, d, g */
        HAL_GPIO_WritePin(
            GPIOB,
            GPIO_PIN_0 | GPIO_PIN_1 | GPIO_PIN_2 |
            GPIO_PIN_3 | GPIO_PIN_6,
            GPIO_PIN_RESET
        );
    }
}
/* USER CODE END 0 */
\end{stmcode}

\subsubsection{Initialization in main}

After GPIO and TIM2 initialization, all four
displays are disabled by setting PA6--PA9 high.
PA4 is also set high so that the DOT LEDs
start in the off state.

The application then prepares digit 1 and enables
the first display by setting PA6 low.
Finally, TIM2 is started in interrupt mode.

\begin{stmcode}{Exercise 2 initialization inside main}
MX_GPIO_Init();
MX_TIM2_Init();

/* USER CODE BEGIN 2 */

/* Disable all four displays */
HAL_GPIO_WritePin(
    GPIOA,
    GPIO_PIN_6 | GPIO_PIN_7 | GPIO_PIN_8 | GPIO_PIN_9,
    GPIO_PIN_SET
);

/* Start with the DOT LEDs off */
HAL_GPIO_WritePin(GPIOA, GPIO_PIN_4, GPIO_PIN_SET);

/* Show 1 on the first display immediately */
display7SEG(1);
HAL_GPIO_WritePin(GPIOA, GPIO_PIN_6, GPIO_PIN_RESET);

/* Start TIM2 and its interrupt */
HAL_TIM_Base_Start_IT(&htim2);

/* USER CODE END 2 */
\end{stmcode}

The main loop contains no application processing.
Digit selection and DOT toggling are both handled
in the timer callback without using
\texttt{HAL\_Delay()}.

\subsubsection{Timer interrupt callback}

The callback checks that the interrupt belongs
to TIM2. Three static variables retain their values
between callback invocations:
\begin{itemize}
    \item \texttt{scan\_ticks} counts timer events
          for the 500-ms display-selection interval.
    \item \texttt{dot\_ticks} counts timer events
          for the one-second DOT-toggle interval.
    \item \texttt{active\_display} identifies the
          selected display using indices 0--3.
\end{itemize}

The initial value of \texttt{active\_display} is 0,
matching the first display enabled during startup.

\begin{stmcode}{Exercise 2 TIM2 callback}
/* USER CODE BEGIN 4 */
void HAL_TIM_PeriodElapsedCallback(TIM_HandleTypeDef *htim)
{
    if (htim->Instance == TIM2)
    {
        static uint8_t scan_ticks = 0;
        static uint8_t dot_ticks = 0;
        static uint8_t active_display = 0;

        /* This callback runs once every 10 ms */
        scan_ticks++;
        dot_ticks++;

        /* 100 * 10 ms = 1 second: change the DOT LEDs */
        if (dot_ticks >= 100)
        {
            dot_ticks = 0;
            HAL_GPIO_TogglePin(GPIOA, GPIO_PIN_4);
        }

        /* 50 * 10 ms = 500 ms: select the next display */
        if (scan_ticks >= 50)
        {
            scan_ticks = 0;

            /* Disable all displays before changing segments */
            HAL_GPIO_WritePin(
                GPIOA,
                GPIO_PIN_6 | GPIO_PIN_7 | GPIO_PIN_8 | GPIO_PIN_9,
                GPIO_PIN_SET
            );

            active_display++;
            if (active_display >= 4)
            {
                active_display = 0;
            }

            if (active_display == 0)
            {
                display7SEG(1);
                HAL_GPIO_WritePin(GPIOA, GPIO_PIN_6, GPIO_PIN_RESET);
            }
            else if (active_display == 1)
            {
                display7SEG(2);
                HAL_GPIO_WritePin(GPIOA, GPIO_PIN_7, GPIO_PIN_RESET);
            }
            else if (active_display == 2)
            {
                display7SEG(3);
                HAL_GPIO_WritePin(GPIOA, GPIO_PIN_8, GPIO_PIN_RESET);
            }
            else
            {
                display7SEG(0);
                HAL_GPIO_WritePin(GPIOA, GPIO_PIN_9, GPIO_PIN_RESET);
            }
        }
    }
}
/* USER CODE END 4 */
\end{stmcode}

After every 50 interrupts, the callback disables
all displays, advances \texttt{active\_display},
writes the required digit pattern, and enables
the corresponding display. When the index reaches
4, it returns to 0. The repeating selection
sequence is therefore
\[
    0 \rightarrow 1 \rightarrow 2 \rightarrow 3
    \rightarrow 0.
\]

Disabling every display before changing the segment
signals prevents the previous display from briefly
showing the next digit's pattern.

The DOT and scanning conditions are separate
\texttt{if} statements. Every 100 interrupts,
both conditions are satisfied in the same callback
invocation: DOT changes state, and the next display
is selected. In the supplied source, DOT is
processed first.

\subsubsection{Expected simulation behavior}

At startup, the first display shows 1 and DOT is off.
Every 500~ms, the next display is selected.
The digits remain fixed at 1, 2, 3, and 0;
this exercise does not advance the displayed time.

\begin{table}[H]
    \centering
    \caption{Nominal display and DOT sequence for Exercise 2}
    \label{tab:ex02-sequence}
    \begin{tabular}{llll}
        \toprule
        \textbf{Time interval} &
        \textbf{Enabled display} &
        \textbf{Digit} &
        \textbf{DOT state} \\
        \midrule
        0--0.5~s   & First, PA6 low  & 1 & Off \\
        0.5--1.0~s & Second, PA7 low & 2 & Off \\
        1.0--1.5~s & Third, PA8 low  & 3 & On \\
        1.5--2.0~s & Fourth, PA9 low & 0 & On \\
        2.0--2.5~s & First, PA6 low  & 1 & Off \\
        \bottomrule
    \end{tabular}
\end{table}

All unselected digit-enable outputs remain high.
The slow 500-ms selection interval makes the
scanning sequence visible in Proteus.

DOT stays off for one second and on for one second.
Therefore, although its state changes every second,
its complete on--off cycle lasts
\[
    T_{\mathrm{DOT\ cycle}} = 2~\mathrm{s},
    \qquad
    f_{\mathrm{DOT\ blink}} = 0.5~\mathrm{Hz}.
\]

\subsubsection{Short question: Frequency of the scanning process}

Following the definition in the introduction,
one complete scan must visit all four displays.

Each display is selected for
\[
    T_S = 500~\mathrm{ms} = 0.5~\mathrm{s}.
\]

Therefore, the complete scanning period is
\[
    T_{\mathrm{scan}}
    = 4T_S
    = 4 \times 0.5~\mathrm{s}
    = 2~\mathrm{s}.
\]

The scanning frequency is
\[
    \boxed{
        f_{\mathrm{scan}}
        = \frac{1}{T_{\mathrm{scan}}}
        = \frac{1}{2~\mathrm{s}}
        = 0.5~\mathrm{Hz}
    }.
\]

The display-selection rate is
\[
    f_{\mathrm{selection}}
    = \frac{1}{T_S}
    = \frac{1}{0.5~\mathrm{s}}
    = 2~\mathrm{Hz}.
\]

Thus, TIM2 interrupts occur at 100~Hz, a new
display is selected twice per second, and the
complete four-display scanning frequency is
\textbf{0.5~Hz}.

% ============================================================
% Replace the existing Exercise 3 subsection with this section.
% The existing reference to the Exercise 2 schematic is preserved.
% ============================================================
\subsection{Exercise 3: Buffered digit selection}

\subsubsection{Objective and circuit}

Exercise 3 introduces an array named \texttt{led\_buffer}
and a function named \texttt{update7SEG(int index)}.
The array stores the digits to display, while the
function selects a physical display and writes
its corresponding array value.

The circuit remains the same as in
Figure~\ref{fig:ex02}. PB0--PB6 carry the shared
segment signals, PA6--PA9 select the four displays,
and PA4 controls the DOT LEDs. These signals
remain active low.

The initial buffer contains 1, 2, 3, and 4.
Display selection remains in the TIM2 callback,
with the next display selected every 500~ms.
DOT continues to change state every second.

\subsubsection{Buffer declarations and function prototypes}

The source declares the following variables in
\texttt{USER CODE BEGIN PV}:
\begin{itemize}
    \item \texttt{MAX\_LED} specifies the number of displays.
    \item \texttt{index\_led} identifies the currently
          selected display, using indices 0--3.
    \item \texttt{led\_buffer} stores the four digit values.
\end{itemize}

\begin{stmcode}{Exercise 3 buffer and function prototypes}
/* USER CODE BEGIN PV */
const int MAX_LED = 4;
int index_led = 0;
int led_buffer[4] = {1, 2, 3, 4};
/* USER CODE END PV */

/* USER CODE BEGIN PFP */
void display7SEG(int num);
void update7SEG(int index);
/* USER CODE END PFP */
\end{stmcode}

The display index and the digit value have different
roles. For example, \texttt{update7SEG(2)} selects
the third physical display and reads
\texttt{led\_buffer[2]}, whose initial value is 3.

\begin{table}[H]
    \centering
    \caption{Initial buffer-to-display mapping for Exercise 3}
    \label{tab:ex03-buffer}
    \begin{tabular}{cccc}
        \toprule
        \textbf{Index} &
        \textbf{Physical display} &
        \textbf{Enable pin} &
        \textbf{Buffer value} \\
        \midrule
        0 & First  & PA6 & 1 \\
        1 & Second & PA7 & 2 \\
        2 & Third  & PA8 & 3 \\
        3 & Fourth & PA9 & 4 \\
        \bottomrule
    \end{tabular}
\end{table}

\subsubsection{Digit encoding for values 0--9}

The Exercise 2 digit function supports only values
0--3. In Exercise 3, \texttt{display7SEG()} is
expanded to support all decimal digits, including
the value 4 required by the initial buffer.

The function builds a pin mask named
\texttt{segments\_on}. It first sets all segment
outputs high, then drives the selected outputs low.
An unsupported value leaves all segments off.

\begin{stmcode}{Exercise 3 display7SEG function}
/* Put a number from 0 to 9 on the shared PB0-PB6 segment lines */
void display7SEG(int num){
    const uint16_t all_segments =
        GPIO_PIN_0 | GPIO_PIN_1 | GPIO_PIN_2 | GPIO_PIN_3 |
        GPIO_PIN_4 | GPIO_PIN_5 | GPIO_PIN_6;

    uint16_t segments_on = 0;

    switch (num)
    {
        case 0: /* a b c d e f */
            segments_on = GPIO_PIN_0 | GPIO_PIN_1 | GPIO_PIN_2 |
                          GPIO_PIN_3 | GPIO_PIN_4 | GPIO_PIN_5;
            break;

        case 1: /* b c */
            segments_on = GPIO_PIN_1 | GPIO_PIN_2;
            break;

        case 2: /* a b d e g */
            segments_on = GPIO_PIN_0 | GPIO_PIN_1 | GPIO_PIN_3 |
                          GPIO_PIN_4 | GPIO_PIN_6;
            break;

        case 3: /* a b c d g */
            segments_on = GPIO_PIN_0 | GPIO_PIN_1 | GPIO_PIN_2 |
                          GPIO_PIN_3 | GPIO_PIN_6;
            break;

        case 4: /* b c f g */
            segments_on = GPIO_PIN_1 | GPIO_PIN_2 |
                          GPIO_PIN_5 | GPIO_PIN_6;
            break;

        case 5: /* a c d f g */
            segments_on = GPIO_PIN_0 | GPIO_PIN_2 | GPIO_PIN_3 |
                          GPIO_PIN_5 | GPIO_PIN_6;
            break;

        case 6: /* a c d e f g */
            segments_on = GPIO_PIN_0 | GPIO_PIN_2 | GPIO_PIN_3 |
                          GPIO_PIN_4 | GPIO_PIN_5 | GPIO_PIN_6;
            break;

        case 7: /* a b c */
            segments_on = GPIO_PIN_0 | GPIO_PIN_1 | GPIO_PIN_2;
            break;

        case 8: /* a b c d e f g */
            segments_on = all_segments;
            break;

        case 9: /* a b c d f g */
            segments_on = GPIO_PIN_0 | GPIO_PIN_1 | GPIO_PIN_2 |
                          GPIO_PIN_3 | GPIO_PIN_5 | GPIO_PIN_6;
            break;

        default:
            break; /* Unsupported number: leave all segments off */
    }

    /* Common anode: HIGH = segment off, LOW = segment on */
    HAL_GPIO_WritePin(GPIOB, all_segments, GPIO_PIN_SET);

    if (segments_on != 0)
    {
        HAL_GPIO_WritePin(GPIOB, segments_on, GPIO_PIN_RESET);
    }
  }
\end{stmcode}

\subsubsection{Report 1: update7SEG function}

The function first checks that the index is between
0 and 3. An invalid index causes it to return
without changing the GPIO outputs.

For a valid index, it disables all four displays
before updating the shared segment signals.
The \texttt{switch} statement then reads the
corresponding buffer value and enables the
selected display.

Both \texttt{display7SEG()} and
\texttt{update7SEG()} are defined inside
\texttt{USER CODE BEGIN 0}.

\begin{stmcode}{Exercise 3 update7SEG function}
/* Select display index 0, 1, 2, or 3 */
void update7SEG(int index){
    if (index < 0 || index >= MAX_LED)
    {
        return;
    }

    /* Disable every display before changing the shared segments */
    HAL_GPIO_WritePin(
        GPIOA,
        GPIO_PIN_6 | GPIO_PIN_7 | GPIO_PIN_8 | GPIO_PIN_9,
        GPIO_PIN_SET
    );

    switch (index)
    {
        case 0:
            display7SEG(led_buffer[0]);
            HAL_GPIO_WritePin(GPIOA, GPIO_PIN_6, GPIO_PIN_RESET);
            break;

        case 1:
            display7SEG(led_buffer[1]);
            HAL_GPIO_WritePin(GPIOA, GPIO_PIN_7, GPIO_PIN_RESET);
            break;

        case 2:
            display7SEG(led_buffer[2]);
            HAL_GPIO_WritePin(GPIOA, GPIO_PIN_8, GPIO_PIN_RESET);
            break;

        case 3:
            display7SEG(led_buffer[3]);
            HAL_GPIO_WritePin(GPIOA, GPIO_PIN_9, GPIO_PIN_RESET);
            break;

        default:
            break;
    }
  }
\end{stmcode}

Disabling all displays before writing the segment
pattern prevents the previously selected display
from briefly showing the next digit.

The function selects and updates one display.
It does not advance \texttt{index\_led} or introduce
a delay; the callback controls the selection order
and timing.

\subsubsection{Initialization in main}

After GPIO and TIM2 initialization, the application
sets DOT high so that its LEDs start off.
It then calls \texttt{update7SEG(index\_led)}.
Since \texttt{index\_led} initially equals 0,
the first display immediately shows
\texttt{led\_buffer[0]}, which is 1.

TIM2 is then started in interrupt mode.

\begin{stmcode}{Exercise 3 initialization inside main}
MX_GPIO_Init();
MX_TIM2_Init();

/* USER CODE BEGIN 2 */
HAL_GPIO_WritePin(GPIOA, GPIO_PIN_4, GPIO_PIN_SET); /* DOT off */
update7SEG(index_led);                              /* Show 1 on display 0 */
HAL_TIM_Base_Start_IT(&htim2);                       /* Start 10 ms interrupts */
/* USER CODE END 2 */
\end{stmcode}

The main loop contains no application processing
in this exercise.

\subsubsection{Report 2: Timer interrupt callback}

The callback retains two static counters.
\texttt{scan\_ticks} controls display selection,
while \texttt{dot\_ticks} controls DOT toggling.
Both counters advance during each 10-ms TIM2 event.

When \texttt{scan\_ticks} reaches 50, the callback
increments \texttt{index\_led}. If the index reaches
\texttt{MAX\_LED}, it returns to 0. The callback
then invokes \texttt{update7SEG(index\_led)}.

When \texttt{dot\_ticks} reaches 100, the callback
toggles PA4. The DOT and scanning conditions are
separate, allowing both tasks to execute during
the same callback invocation when required.

\begin{stmcode}{Exercise 3 TIM2 callback}
/* USER CODE BEGIN 4 */
void HAL_TIM_PeriodElapsedCallback(TIM_HandleTypeDef *htim)
{
    if (htim->Instance == TIM2)
    {
        static uint8_t scan_ticks = 0;
        static uint8_t dot_ticks = 0;

        scan_ticks++;
        dot_ticks++;

        /* Change the DOT LEDs every 100 * 10 ms = 1 second */
        if (dot_ticks >= 100)
        {
            dot_ticks = 0;
            HAL_GPIO_TogglePin(GPIOA, GPIO_PIN_4);
        }

        /* Select the next display every 50 * 10 ms = 500 ms */
        if (scan_ticks >= 50)
        {
            scan_ticks = 0;

            index_led++;
            if (index_led >= MAX_LED)
            {
                index_led = 0;
            }

            update7SEG(index_led);
        }
    }
}
/* USER CODE END 4 */
\end{stmcode}

Unlike Exercise 2, the callback no longer contains
the individual digit patterns and enable-pin
selections. Those operations are handled by
\texttt{update7SEG()}, while the displayed values
are stored in \texttt{led\_buffer}.

\subsubsection{Expected behavior and buffer testing}

With the initial buffer \texttt{\{1, 2, 3, 4\}},
the displays are selected in physical order:
first, second, third, and fourth. Each remains
selected for 500~ms before the next display is
enabled. DOT changes state every second.

To test the buffer-based selection, change the
initializer of \texttt{led\_buffer}, rebuild the
program, and run the updated firmware in Proteus.
The following cases provide expected results.

\begin{table}[H]
    \centering
    \caption{Suggested buffer tests for Exercise 3}
    \label{tab:ex03-tests}
    \begin{tabular}{ll}
        \toprule
        \textbf{Buffer initializer} &
        \textbf{Expected digits, first to fourth} \\
        \midrule
        \texttt{\{1, 2, 3, 4\}} & 1, 2, 3, 4 \\
        \texttt{\{4, 3, 2, 1\}} & 4, 3, 2, 1 \\
        \texttt{\{9, 0, 5, 8\}} & 9, 0, 5, 8 \\
        \bottomrule
    \end{tabular}
\end{table}

Changing the buffer changes the digit assigned to
each display. The physical selection order,
selection interval, and DOT timing remain the same.
% ============================================================
% ============================================================
% Replace the existing Exercise 4 subsection with this section.
% ============================================================
\subsection{Exercise 4: One complete scan per second}

\subsubsection{Objective and circuit}

Exercise 4 changes the display-selection interval
so that one complete scan of all four seven-segment
displays occurs every second. The required complete
scanning frequency is therefore 1~Hz.
DOT continues to change state every second.

The circuit remains the same as in
Figure~\ref{fig:ex02}. The GPIO assignments and
the implementations of \texttt{display7SEG()}
and \texttt{update7SEG()} are unchanged from
Exercise 3.

The source retains the following buffer declarations.

\begin{stmcode}{Exercise 4 display buffer}
/* USER CODE BEGIN PV */
const int MAX_LED = 4;
int index_led = 0;
int led_buffer[4] = {1, 2, 3, 4};
/* USER CODE END PV */
\end{stmcode}

The displayed values remain 1, 2, 3, and 4.
The change from Exercise 3 is the scanning interval:
each display is now selected for 250~ms instead
of 500~ms.

\subsubsection{Required scan timing}

TIM2 retains its 10-ms interrupt interval.
The following calculation determines how often
the callback must select the next display.

\paragraph{Step 1: Determine the complete scanning period.}

The required scanning frequency is
\[
    f_{\mathrm{scan}} = 1~\mathrm{Hz}.
\]

Therefore,
\[
    T_{\mathrm{scan}}
    = \frac{1}{f_{\mathrm{scan}}}
    = 1~\mathrm{s}
    = 1000~\mathrm{ms}.
\]

\paragraph{Step 2: Determine the interval for one display.}

One complete scan contains four display-selection
intervals. Each display must therefore remain
selected for
\[
    T_S
    = \frac{T_{\mathrm{scan}}}{4}
    = \frac{1000~\mathrm{ms}}{4}
    = 250~\mathrm{ms}.
\]

\paragraph{Step 3: Determine the callback counter threshold.}

Since each TIM2 event occurs every 10~ms,
the required number of interrupts is
\[
    N_{\mathrm{scan}}
    = \frac{250~\mathrm{ms}}{10~\mathrm{ms}}
    = 25.
\]

Thus, the scanning condition changes from
\texttt{scan\_ticks >= 50} in Exercise 3 to
\texttt{scan\_ticks >= 25} in Exercise 4.

\paragraph{Step 4: Verify the complete scanning frequency.}

The four display intervals give
\[
    T_{\mathrm{scan}}
    = 4 \times 250~\mathrm{ms}
    = 1000~\mathrm{ms}
    = 1~\mathrm{s}.
\]

Therefore,
\[
    \boxed{
        f_{\mathrm{scan}}
        = \frac{1}{1~\mathrm{s}}
        = 1~\mathrm{Hz}
    }.
\]

The callback selects a new display four times
per second. Each individual display is visited
once per second.

\subsubsection{Initialization in main}

Initialization remains the same as in Exercise 3.
DOT starts off, the first display immediately
shows \texttt{led\_buffer[0]}, and TIM2 is
started in interrupt mode.

\begin{stmcode}{Exercise 4 initialization inside main}
MX_GPIO_Init();
MX_TIM2_Init();

/* USER CODE BEGIN 2 */
HAL_GPIO_WritePin(GPIOA, GPIO_PIN_4, GPIO_PIN_SET); /* DOT off */
update7SEG(index_led);                              /* Show 1 on display 0 */
HAL_TIM_Base_Start_IT(&htim2);                       /* Start 10 ms interrupts */
/* USER CODE END 2 */
\end{stmcode}

The main loop contains no application processing.
Display selection and DOT toggling remain in
the timer callback.

\subsubsection{Report 1: Timer interrupt callback}

The callback uses two static counters:
\texttt{scan\_ticks} for display selection and
\texttt{dot\_ticks} for DOT toggling. Both
counters advance on each TIM2 event.

The DOT threshold remains 100 interrupts.
The scanning threshold is reduced to 25 interrupts.

\begin{stmcode}{Exercise 4 TIM2 callback}
/* USER CODE BEGIN 4 */
void HAL_TIM_PeriodElapsedCallback(TIM_HandleTypeDef *htim)
{
    if (htim->Instance == TIM2)
    {
        static uint8_t scan_ticks = 0;
        static uint8_t dot_ticks = 0;

        scan_ticks++;
        dot_ticks++;

        /* Change the DOT LEDs every 100 * 10 ms = 1 second */
        if (dot_ticks >= 100)
        {
            dot_ticks = 0;
            HAL_GPIO_TogglePin(GPIOA, GPIO_PIN_4);
        }

        /* Select the next display every 25 * 10 ms = 250 ms */
        if (scan_ticks >= 25)
        {
            scan_ticks = 0;

            index_led++;
            if (index_led >= MAX_LED)
            {
                index_led = 0;
            }

            update7SEG(index_led);
        }
    }
}
/* USER CODE END 4 */
\end{stmcode}

When \texttt{scan\_ticks} reaches 25, it resets
to zero. The callback advances \texttt{index\_led},
returns it to zero when it reaches 4, and calls
\texttt{update7SEG(index\_led)}.

The selected indices repeat in the order
\[
    0 \rightarrow 1 \rightarrow 2 \rightarrow 3
    \rightarrow 0.
\]

The \texttt{update7SEG()} function disables all
four displays before changing the shared segment
pattern, then enables the selected display.

The DOT and scanning conditions are separate
\texttt{if} statements. Every 100 interrupts,
both conditions are satisfied in the same callback
invocation. The source toggles DOT first and
then selects the next display.

\subsubsection{Expected simulation behavior}

The first display shows 1 immediately after
initialization. Every 250~ms, the next display
is selected. After one second, the selection
returns to the first display.

\begin{table}[H]
    \centering
    \caption{Nominal display sequence for Exercise 4}
    \label{tab:ex04-sequence}
    \begin{tabular}{llll}
        \toprule
        \textbf{Time interval} &
        \textbf{Enabled display} &
        \textbf{Digit} &
        \textbf{DOT state} \\
        \midrule
        0--0.25~s    & First, PA6 low  & 1 & Off \\
        0.25--0.50~s & Second, PA7 low & 2 & Off \\
        0.50--0.75~s & Third, PA8 low  & 3 & Off \\
        0.75--1.00~s & Fourth, PA9 low & 4 & Off \\
        1.00--1.25~s & First, PA6 low  & 1 & On \\
        \bottomrule
    \end{tabular}
\end{table}

Only the selected display is enabled; the other
digit-enable outputs remain high. DOT remains
off for one second and on for one second,
giving a complete blinking cycle of two seconds.

Compared with Exercise 3, the complete scanning
frequency increases from 0.5~Hz to 1~Hz.
The TIM2 interrupt interval and DOT timing
remain unchanged.

The required 1-Hz scan is intentionally slow,
so the alternating displays remain visible
in Proteus.
% ============================================================
% ============================================================
% Replace the existing Exercise 5 subsection with this section.
% ============================================================
\subsection{Exercise 5: Digital clock HH:MM}

\subsubsection{Objective and inherited configuration}

Exercise 5 displays a 24-hour digital clock.
The first two displays show the hour, and the
last two show the minute. Seconds are maintained
in software but are not displayed.

The circuit remains as shown in Figure~\ref{fig:ex02}.
GPIO configuration, TIM2 configuration,
\texttt{display7SEG()}, \texttt{update7SEG()},
and the timer callback are already established
in Exercise 4 and are retained.

The callback continues selecting a display every
250~ms and toggling DOT every second.
The new clock processing is performed in
\texttt{main()}.

\subsubsection{Time variables and display buffer}

The previous buffer declaration is replaced with
a volatile buffer initialized to zero, and three
time variables are added. The declarations of
\texttt{MAX\_LED} and \texttt{index\_led} remain
as in Exercise 4.

\begin{stmcode}{Exercise 5 updated buffer and new time variables}
volatile int led_buffer[4] = {0, 0, 0, 0};

int hour = 15;
int minute = 8;
int second = 50;
\end{stmcode}

These declarations belong in
\texttt{USER CODE BEGIN PV}.
The initial time is 15:08:50.

The buffer is written by the main loop and read
during interrupt-driven display selection.
Its declaration uses \texttt{volatile} as in
the supplied source.

The new function prototype
\texttt{void updateClockBuffer(void);}
is added to \texttt{USER CODE BEGIN PFP},
alongside the existing prototypes.

\subsubsection{Report 1: updateClockBuffer function}

The function converts the hour and minute into
four decimal digits. Integer division by 10
extracts the tens digit, while the remainder
after division by 10 extracts the units digit.

\begin{stmcode}{Exercise 5 updateClockBuffer function}
void updateClockBuffer(void){
    led_buffer[0] = hour / 10;    /* Hour tens */
    led_buffer[1] = hour % 10;    /* Hour ones */
    led_buffer[2] = minute / 10;  /* Minute tens */
    led_buffer[3] = minute % 10;  /* Minute ones */
}
\end{stmcode}

The definition is placed in
\texttt{USER CODE BEGIN 0}.

For the initial time 15:08, the function produces
\texttt{\{1, 5, 0, 8\}}. A single-digit hour
or minute automatically receives a leading zero.
For example, 07:05 produces
\texttt{\{0, 7, 0, 5\}}.

\subsubsection{Initialization}

Initialization follows Exercise 4, with one added
call to \texttt{updateClockBuffer()} before
\texttt{update7SEG(index\_led)}.

This fills the initially zeroed buffer with
\texttt{\{1, 5, 0, 8\}} before the first display
is enabled and TIM2 is started. DOT still
starts in the off state.

\subsubsection{Clock update in the main loop}

The previously empty main loop now waits
approximately one second, increments the seconds,
handles time rollover, and updates the buffer.

The supplied source places \texttt{HAL\_Delay(1000)}
before \texttt{second++}, so the initial time
is retained during the first waiting interval.

\begin{stmcode}{Exercise 5 main loop}
/* USER CODE BEGIN WHILE */
while (1) {
    HAL_Delay(1000); /* Wait approximately 1 second */
    second++;

    if (second >= 60) {
      second = 0;
      minute++;
    }

    if (minute >= 60) {
      minute = 0;
      hour++;
    }

    if (hour >= 24) {
      hour = 0;
    }

    updateClockBuffer();
    /* USER CODE END WHILE */

    /* USER CODE BEGIN 3 */
}
/* USER CODE END 3 */
\end{stmcode}

When seconds reach 60, they reset to zero and
the minute increases. When minutes reach 60,
they reset to zero and the hour increases.
When hours reach 24, they reset to zero.

The separate conditions allow all required
rollovers to occur during the same iteration.
For example, 23:59:59 advances to 00:00:00.

Although \texttt{HAL\_Delay()} blocks progress
in the main loop, TIM2 interrupts remain enabled.
The callback therefore continues scanning the
displays and toggling DOT during the delay.

\subsubsection{Expected behavior}

The clock initially displays 15:08, with
50 seconds stored internally. After approximately
ten seconds, the internal time reaches 15:09:00
and the displayed digits change to 15:09.

The hour and minute digits change when their
corresponding values advance. Display scanning
and DOT blinking continue throughout the clock
update process.

% ============================================================
% ============================================================
% Replace the existing Exercise 6 subsection with this section.
% Only new or changed code is presented.
% ============================================================
% ============================================================
% Replace BOTH existing Exercise 6 and Exercise 7 sections
% with this complete block.
% Stop replacing immediately before the Exercise 8 heading.
% ============================================================

\subsection{Exercise 6: Software timer for the red LED}

\subsubsection{Objective and inherited configuration}

Exercise 6 introduces a software timer to control
the red LED on PA5 and removes the blocking
\texttt{HAL\_Delay()} from the clock-processing loop.

The circuit remains as shown in Figure~\ref{fig:ex02}.
GPIO and TIM2 configuration, the display buffer,
and the functions \texttt{display7SEG()},
\texttt{update7SEG()}, and
\texttt{updateClockBuffer()} are retained
from Exercise 5.

TIM2 continues generating interrupts every 10~ms.
The callback advances the software timer, while
the main loop handles the LED and clock events.

\subsubsection{New variables and software-timer functions}

The following declarations are added to
\texttt{USER CODE BEGIN PV}.

\begin{stmcode}{Exercise 6 new variables}
volatile int timer0_counter = 0;
volatile int timer0_flag = 0;
const int TIMER_CYCLE = 10;       /* TIM2 period: 10 ms */
volatile int clock_second_flag = 0;
\end{stmcode}

\texttt{timer0\_counter} stores the remaining
number of 10-ms timer events.
\texttt{timer0\_flag} indicates expiry.
The separate \texttt{clock\_second\_flag}
notifies the main loop that one second has passed.

The following functions are added to
\texttt{USER CODE BEGIN 0}.

\begin{stmcode}{Exercise 6 software-timer functions}
void setTimer0(int duration)
{
    timer0_counter = duration / TIMER_CYCLE;
    timer0_flag = 0;
}

void timer_run(void)
{
    if (timer0_counter > 0)
    {
        timer0_counter--;

        if (timer0_counter == 0)
        {
            timer0_flag = 1;
        }
    }
}
\end{stmcode}

\texttt{setTimer0()} converts a duration in
milliseconds into timer events and clears the
previous expiry flag.

\texttt{timer\_run()} decreases a positive
counter once per callback. When it reaches
zero, the function raises the flag.
The main loop performs the LED action.

\subsubsection{Updated timer interrupt callback}

The function
\texttt{HAL\_TIM\_PeriodElapsedCallback()}
is modified to call \texttt{timer\_run()}
and raise \texttt{clock\_second\_flag}.

Display selection and DOT toggling remain
inside the callback.

\begin{stmcode}{Exercise 6 updated timer callback}
/* USER CODE BEGIN 4 */
void HAL_TIM_PeriodElapsedCallback(TIM_HandleTypeDef *htim)
{
    if (htim->Instance == TIM2)
    {
        static uint8_t scan_ticks = 0;
        static uint8_t dot_ticks = 0;

        timer_run();    /* Decrease software timer once per 10 ms */

        scan_ticks++;
        dot_ticks++;

        if (dot_ticks >= 100)
        {
            dot_ticks = 0;
            HAL_GPIO_TogglePin(GPIOA, GPIO_PIN_4);
            clock_second_flag = 1; /* Tell main one second passed */
        }

        if (scan_ticks >= 25)
        {
            scan_ticks = 0;

            index_led++;
            if (index_led >= MAX_LED)
            {
                index_led = 0;
            }

            update7SEG(index_led);
        }
    }
}
/* USER CODE END 4 */
\end{stmcode}

The timer-instance check restricts processing
to TIM2. Each callback advances timer 0.
Every 100 callbacks, DOT changes state and
the clock flag is raised. Every 25 callbacks,
the next display is selected.

Thus, the clock notification is generated
by TIM2, while the main loop performs the
time update.

\subsubsection{Initialization}

The display initialization follows Exercise 5.
PA5 is also set high so that the red LED starts
off, and timer 0 is initialized before TIM2 starts.

\begin{stmcode}{Exercise 6 initialization}
/* USER CODE BEGIN 2 */
HAL_GPIO_WritePin(
    GPIOA, GPIO_PIN_4 | GPIO_PIN_5, GPIO_PIN_SET
);                         /* DOT and PA5 LED initially off */

updateClockBuffer();
update7SEG(index_led);      /* Display 0 initially selected */

setTimer0(1000);            /* First PA5 toggle after 1 second */
HAL_TIM_Base_Start_IT(&htim2);
/* USER CODE END 2 */
\end{stmcode}

\subsubsection{Updated main loop}

The main loop checks two independent flags.
Timer 0 controls the PA5 LED, while
\texttt{clock\_second\_flag} controls
the clock update.

\begin{stmcode}{Exercise 6 main loop}
/* USER CODE BEGIN WHILE */
while (1) {
  if (timer0_flag == 1) {
    HAL_GPIO_TogglePin(GPIOA, GPIO_PIN_5);
    setTimer0(2000);
  }

  if (clock_second_flag == 1) {
    clock_second_flag = 0;

    second++;

    if (second >= 60) {
      second = 0;
      minute++;
    }

    if (minute >= 60) {
      minute = 0;
      hour++;
    }

    if (hour >= 24) {
      hour = 0;
    }

    updateClockBuffer();
  }
    /* USER CODE END WHILE */

    /* USER CODE BEGIN 3 */
}
/* USER CODE END 3 */
\end{stmcode}

The first LED toggle occurs after approximately
one second. Subsequent toggles occur every
two seconds because the main loop rearms
timer 0 with \texttt{setTimer0(2000)}.
This call also clears its expiry flag.

The clock-rollover logic follows Exercise 5.
It now executes when the clock flag is raised,
without a blocking delay.

\subsubsection{Answers to the three report questions}

Each question changes only the initial
\texttt{setTimer0(1000)} call.
The later \texttt{setTimer0(2000)} call
remains unchanged.

\begin{table}[H]
    \centering
    \caption{Effect of changing the initial software-timer call}
    \label{tab:ex06-questions}
    \begin{tabular}{
        >{\raggedright\arraybackslash}p{0.28\textwidth}
        >{\raggedright\arraybackslash}p{0.61\textwidth}}
        \toprule
        \textbf{Initial call} &
        \textbf{Outcome and reason} \\
        \midrule

        Omitted &
        The counter and flag remain zero.
        Since \texttt{timer\_run()} only processes
        positive counters, no expiry flag is raised.
        The PA5 LED remains off. \\

        \addlinespace
        \texttt{setTimer0(1)} &
        Integer division gives
        \texttt{1 / 10 = 0}, and the flag is cleared.
        The zero counter is not decremented,
        so the flag is never raised.
        The PA5 LED remains off. \\

        \addlinespace
        \texttt{setTimer0(10)} &
        Integer division gives
        \texttt{10 / 10 = 1}.
        The next TIM2 callback decrements the
        counter to zero and raises the flag.
        The first toggle occurs after approximately
        10~ms. Subsequent toggles occur every
        two seconds after rearming. \\

        \bottomrule
    \end{tabular}
\end{table}

These changes affect the PA5 task.
Display scanning, DOT blinking, and clock
updates continue independently of timer 0.

\subsubsection{Expected behavior}

PA5 initially remains off, turns on after
approximately one second, and then changes
state every two seconds.

The initial clock time is 15:08:50.
The displayed time becomes 15:09 after
approximately ten seconds. Display selection
continues every 250~ms, and DOT changes
state every second.


% ============================================================
\subsection{Exercise 7: Clock and DOT handled in main}

\subsubsection{Objective and inherited configuration}

Exercise 7 uses a software timer for the clock
and moves DOT toggling to the main loop.
The circuit remains as shown in
Figure~\ref{fig:ex02}.

GPIO and TIM2 configuration, display functions,
clock-buffer conversion, and
\texttt{setTimer0()} are already established
in the previous exercises and are retained.

In the supplied source, timer 0 is reassigned
to the one-second clock and DOT task.
The previous \texttt{clock\_second\_flag}
is removed.

\subsubsection{Software-timer extension}

The source adds the following declarations
to \texttt{USER CODE BEGIN PV}.

\begin{stmcode}{Exercise 7 additional timer variables}
volatile int timer1_counter = 0;
volatile int timer1_flag = 0;
\end{stmcode}

The existing timer 0 declarations and
\texttt{TIMER\_CYCLE} remain unchanged.

A second timer-setting function is added,
and \texttt{timer\_run()} is extended to
process both counters.

\begin{stmcode}{Exercise 7 additional timer and updated timer_run}
void setTimer1(int duration)
{
    timer1_counter = duration / TIMER_CYCLE;
    timer1_flag = 0;
}

void timer_run(void)
{
    if (timer0_counter > 0)
    {
        timer0_counter--;

        if (timer0_counter == 0)
        {
            timer0_flag = 1;
        }
    }

    if (timer1_counter > 0)
    {
        timer1_counter--;

        if (timer1_counter == 0)
        {
            timer1_flag = 1;
        }
    }
}
\end{stmcode}

These functions are defined in
\texttt{USER CODE BEGIN 0}.
Both timers follow the countdown mechanism
established in Exercise 6.

\subsubsection{Updated timer interrupt callback}

The function
\texttt{HAL\_TIM\_PeriodElapsedCallback()}
is modified again. The DOT counter and
DOT-toggle block are removed.

The callback now advances the software timers
and maintains the existing display scan.

\begin{stmcode}{Exercise 7 updated timer callback}
/* USER CODE BEGIN 4 */
void HAL_TIM_PeriodElapsedCallback(TIM_HandleTypeDef *htim)
{
    if (htim->Instance == TIM2)
    {
        static uint8_t scan_ticks = 0;

        timer_run();    /* Count down timer0 and timer1 every 10 ms */

        scan_ticks++;
        if (scan_ticks >= 25)
        {
            scan_ticks = 0;

            index_led++;
            if (index_led >= MAX_LED)
            {
                index_led = 0;
            }

            update7SEG(index_led);
        }
    }
}
/* USER CODE END 4 */
\end{stmcode}

Display selection remains inside the callback
and occurs every 250~ms.
Clock updates and DOT toggling are performed
by the main loop when timer 0 expires.

\subsubsection{Initialization}

Initialization follows Exercise 6, with both
software timers initially set to 1000~ms.

\begin{stmcode}{Exercise 7 initialization}
/* USER CODE BEGIN 2 */
HAL_GPIO_WritePin(
    GPIOA, GPIO_PIN_4 | GPIO_PIN_5, GPIO_PIN_SET
);                          /* DOT and red LED initially off */

updateClockBuffer();
update7SEG(index_led);       /* Select the first clock digit */

setTimer0(1000);             /* Clock and DOT timer */
setTimer1(1000);             /* Initialized but not handled in main */

HAL_TIM_Base_Start_IT(&htim2);
/* USER CODE END 2 */
\end{stmcode}

\subsubsection{Report 1: Main-loop source code}

The main loop checks \texttt{timer0\_flag}.
When timer 0 expires, it rearms the timer,
advances the clock, updates the display buffer,
and toggles DOT.

\begin{stmcode}{Exercise 7 main loop}
/* USER CODE BEGIN WHILE */
while (1) {
    /* USER CODE END WHILE */

    /* USER CODE BEGIN 3 */
    if (timer0_flag == 1) {
      /* Clear the flag and schedule the next second. */
      setTimer0(1000);

      second++;

      if (second >= 60) {
        second = 0;
        minute++;
      }

      if (minute >= 60) {
        minute = 0;
        hour++;
      }

      if (hour >= 24) {
        hour = 0;
      }

      updateClockBuffer();
      HAL_GPIO_TogglePin(GPIOA, GPIO_PIN_4);
    }
  } /* Closes while (1). */
/* USER CODE END 3 */
\end{stmcode}

Calling \texttt{setTimer0(1000)} clears the
expiry flag and schedules the next clock event.
The rollover logic is retained from Exercise 5.

DOT now changes state during the same
main-loop iteration as the clock update.
No \texttt{HAL\_Delay()} is required.

\subsubsection{Expected behavior and timer usage}

The clock starts at 15:08:50 and advances
approximately once per second. DOT also
changes state once per second.
The displayed time becomes 15:09 after
approximately ten seconds.

In the supplied source, the main loop checks
only \texttt{timer0\_flag}.
Timer 1 expires once, and its flag remains
set because no application task handles it.

The main loop contains no PA5 toggle operation.
The red LED therefore remains off, while
timer 0 controls the clock and DOT.
% ============================================================
% Replace the existing Exercise 8 subsection with this section.
% ============================================================
\subsection{Exercise 8: Software-timer-only TIM2 callback}

\subsubsection{Objective and inherited configuration}

Exercise 8 moves periodic display selection
from the timer callback to the main loop.
The interrupt callback now only advances
the software timers.

The circuit remains as shown in Figure~\ref{fig:ex02}.
GPIO and TIM2 configuration match Exercise 7.

The functions \texttt{display7SEG()},
\texttt{update7SEG()}, and
\texttt{updateClockBuffer()} match exactly
the same functions in Exercise 7.
The functions \texttt{setTimer0()} and
\texttt{setTimer1()} also retain the same logic.

The supplied source assigns the three software
timers as follows.

\begin{table}[H]
    \centering
    \caption{Software-timer assignments for Exercise 8}
    \label{tab:ex08-timers}
    \begin{tabular}{
        p{0.13\textwidth}
        p{0.31\textwidth}
        p{0.44\textwidth}}
        \toprule
        \textbf{Timer} &
        \textbf{Task} &
        \textbf{Duration} \\
        \midrule
        Timer 0 & PA5 red LED &
        1000~ms initially; 2000~ms thereafter \\
        Timer 1 & Clock update and DOT &
        1000~ms \\
        Timer 2 & Display selection &
        250~ms \\
        \bottomrule
    \end{tabular}
\end{table}

Unlike the supplied Exercise 7 source,
the clock task now uses timer 1, and the
PA5 red-LED task is active again.

\subsubsection{Additional software timer}

The following declarations are added to
\texttt{USER CODE BEGIN PV}.
All other variable declarations retain
their values from Exercise 7.

\begin{stmcode}{Exercise 8 additional timer variables}
volatile int timer2_counter = 0;
volatile int timer2_flag = 0;
\end{stmcode}

The new function \texttt{setTimer2()} is
defined in \texttt{USER CODE BEGIN 0}.
It follows the same mechanism as
\texttt{setTimer0()} and \texttt{setTimer1()}
in Exercise 7.

\begin{stmcode}{Exercise 8 setTimer2 function}
void setTimer2(int duration) {
  timer2_counter = duration / TIMER_CYCLE;
  timer2_flag = 0;
}
\end{stmcode}

The timer 0 and timer 1 countdown blocks
inside \texttt{timer\_run()} retain the
same logic as Exercise 7. The following
block is added after them, before the
function's closing brace.

\begin{stmcode}{Additional timer 2 block inside timer_run}
if (timer2_counter > 0) {
  timer2_counter--;
  if (timer2_counter == 0) {
    timer2_flag = 1;
  }
}
\end{stmcode}

Thus, each call to \texttt{timer\_run()}
advances all three software timers and
raises the corresponding expiry flags.

\subsubsection{Updated timer interrupt callback}

The function
\texttt{HAL\_TIM\_PeriodElapsedCallback()}
is changed from Exercise 7.
The \texttt{scan\_ticks} counter and
display-selection block are removed.

\begin{stmcode}{Exercise 8 updated timer callback}
/* USER CODE BEGIN 4 */
void HAL_TIM_PeriodElapsedCallback(TIM_HandleTypeDef *htim)
{
    if (htim->Instance == TIM2)
    {
        timer_run();
    }
}
/* USER CODE END 4 */
\end{stmcode}

The callback checks the timer instance and
advances the software timers every 10~ms.
LED toggling, clock updates, and periodic
display selection are handled in the main loop.

\subsubsection{Initialization}

Initialization matches Exercise 7:
DOT and PA5 start off, the clock buffer is
prepared, the first display is selected,
and timers 0 and 1 are initialized to 1000~ms.

One additional call is inserted after
\texttt{setTimer1(1000)} and before
\texttt{HAL\_TIM\_Base\_Start\_IT(\&htim2)}.

\begin{stmcode}{Exercise 8 additional initialization}
setTimer2(250);  /* Next display selection */
\end{stmcode}

\subsubsection{Report 1: Main-loop source code}

The main loop now checks three independent
expiry flags. The clock-rollover logic
matches Exercise 7, but it is triggered
by \texttt{timer1\_flag}.

\begin{stmcode}{Exercise 8 main loop}
/* USER CODE BEGIN WHILE */
while (1) {
    /* Red LED: first toggle after 1 s, then every 2 s */
    if (timer0_flag == 1) {
      HAL_GPIO_TogglePin(GPIOA, GPIO_PIN_5);
      setTimer0(2000);
    }

    /* Digital clock and DOT: update every 1 s */
    if (timer1_flag == 1) {
      setTimer1(1000);

      second++;

      if (second >= 60) {
        second = 0;
        minute++;
      }

      if (minute >= 60) {
        minute = 0;
        hour++;
      }

      if (hour >= 24) {
        hour = 0;
      }

      updateClockBuffer();
      HAL_GPIO_TogglePin(GPIOA, GPIO_PIN_4);
    }

    /* Seven-segment scan: change display every 250 ms */
    if (timer2_flag == 1) {
      setTimer2(250);

      index_led++;
      if (index_led >= MAX_LED) {
        index_led = 0;
      }

      update7SEG(index_led);
    }
    /* USER CODE END WHILE */

    /* USER CODE BEGIN 3 */
}
/* USER CODE END 3 */
\end{stmcode}

Each timer-setting call clears its expiry
flag and schedules the next event.

Timer 2 advances \texttt{index\_led}
through indices 0--3 and calls
\texttt{update7SEG(index\_led)}.
The display-selection function itself
matches Exercise 7; its periodic invocation
has moved to the main loop.

The three conditions are separate
\texttt{if} statements, allowing multiple
expired tasks to be handled during the
same loop iteration.

\subsubsection{Expected behavior}

The clock initially displays 15:08 and
advances approximately once per second.
DOT changes state with each clock update.

The selected display changes every 250~ms.
PA5 starts off, toggles after approximately
one second, and then changes state every
two seconds.

All tasks run without \texttt{HAL\_Delay()},
while the TIM2 callback only maintains
the software-timer countdowns.


% ============================================================
% Replace the existing Exercise 9 and Exercise 10 sections
% with this complete block. Keep the final \end{document}.
%
% Also update the matrix-control GPIO entry in the common
% hardware table to: PA2, PA3, PA10--PA15.
% ============================================================

% ============================================================
% Replace the existing Exercise 9 and Exercise 10 sections
% with this complete block. Keep the final \end{document}.
%
% Updated matrix-control pins: PA2, PA3, PA10--PA15.
% ============================================================

\subsection{Exercise 9: Display the letter A on an
8\texorpdfstring{$\times$}{x}8 matrix}

\subsubsection{Objective and inherited configuration}

Exercise 9 adds an $8 \times 8$ LED matrix and
a ULN2803 column-control stage to display
a stationary capital A.

The functions \texttt{display7SEG()},
\texttt{update7SEG()}, \texttt{updateClockBuffer()},
\texttt{setTimer0()}, \texttt{setTimer1()},
\texttt{setTimer2()}, and \texttt{timer\_run()}
match exactly the same functions in Exercise 8.

TIM2 configuration and
\texttt{HAL\_TIM\_PeriodElapsedCallback()}
also match Exercise 8. TIM2 continues interrupting
every 10~ms, and its callback only calls
\texttt{timer\_run()}.

Matrix scanning is added to the main loop,
alongside the existing clock and LED tasks.

\subsubsection{Updated matrix GPIO configuration}

The updated circuit connects ENM0 to PA2
and ENM1 to PA3. The remaining column
assignments are shown in
Table~\ref{tab:ex09-matrix-pins}.

\begin{table}[H]
    \centering
    \caption{Updated GPIO assignments for the LED matrix}
    \label{tab:ex09-matrix-pins}
    \begin{tabular}{ccc}
        \toprule
        \textbf{Column index} &
        \textbf{Enable signal} &
        \textbf{GPIO pin} \\
        \midrule
        0 & ENM0 & PA2  \\
        1 & ENM1 & PA3  \\
        2 & ENM2 & PA10 \\
        3 & ENM3 & PA11 \\
        4 & ENM4 & PA12 \\
        5 & ENM5 & PA13 \\
        6 & ENM6 & PA14 \\
        7 & ENM7 & PA15 \\
        \bottomrule
    \end{tabular}
\end{table}

All column-control signals use GPIOA.
ROW0--ROW7 connect to PB8--PB15, respectively.

The updated definitions in \texttt{main.h}
include:

\begin{stmcode}{Updated ENM0 and ENM1 pin definitions}
#define ENM0_Pin GPIO_PIN_2
#define ENM0_GPIO_Port GPIOA
#define ENM1_Pin GPIO_PIN_3
#define ENM1_GPIO_Port GPIOA
\end{stmcode}

The generated \texttt{MX\_GPIO\_Init()} function
uses these pin definitions to configure
the matrix outputs as push-pull, with
no internal pull resistors and an initial
HIGH output level.

The matrix-display function also uses
the ENM pin definitions. GPIO initialization
and column selection therefore follow
the same mapping.

In the implemented circuit, HIGH on an
ENM output disables its column, while LOW
selects it. For the selected column,
a LOW row output lights the corresponding
pixel and a HIGH row output turns it off.

\subsubsection{Report 1: Schematic}

\begin{figure}[H]
    \centering
    \schematicplaceholder{figures/pic03.png}{0.92}
    \caption{Exercise 9 Proteus circuit with the
    8$\times$8 LED matrix and ULN2803.}
    \label{fig:ex09}
\end{figure}

\subsubsection{Matrix variables and character pattern}

The following declarations are added inside
\texttt{USER CODE BEGIN PV}.
The existing clock, display-buffer, and
software-timer variables match Exercise 8.

\begin{stmcode}{Exercise 9 matrix variables and smaller A}
const int MAX_LED_MATRIX = 8;
int index_led_matrix = 0;

/* Smaller A: 5 columns wide, 6 rows tall, with blank margins */
uint8_t matrix_buffer[8] = {
    0x00, 0x7C, 0x12, 0x12,
    0x12, 0x7C, 0x00, 0x00
};

uint32_t last_matrix_scan = 0;
\end{stmcode}

Each buffer element represents one column.
Its eight bits specify the pixels on
ROW0--ROW7. Bit 0 corresponds to PB8,
and bit 7 corresponds to PB15.

A stored bit value of 1 represents a lit
pixel and produces a LOW row output.
A stored bit value of 0 produces a HIGH
row output.

The A occupies five columns and six rows.
The remaining positions provide blank margins.

\subsubsection{Report 2: updateLEDMatrix function}

The function \texttt{updateLEDMatrix()}
is defined inside \texttt{USER CODE BEGIN 0}.
It checks the index, disables all columns,
writes the row pattern, and enables the
selected column.

\begin{stmcode}{Exercise 9 updateLEDMatrix function}
void updateLEDMatrix(int index) {
  const uint16_t all_columns = ENM0_Pin | ENM1_Pin | ENM2_Pin | ENM3_Pin |
                               ENM4_Pin | ENM5_Pin | ENM6_Pin | ENM7_Pin;

  uint16_t selected_column;

  if (index < 0 || index >= MAX_LED_MATRIX) {
    return;
  }

  /* Disable all columns before changing the row data. */
  HAL_GPIO_WritePin(GPIOA, all_columns, GPIO_PIN_SET);

  /* Update the eight active-low row outputs. */
  for (int row = 0; row < 8; row++) {
    uint16_t row_pin = (uint16_t)(GPIO_PIN_8 << row);

    GPIO_PinState row_state = (matrix_buffer[index] & (1U << row))
                                  ? GPIO_PIN_RESET /* Pixel ON */
                                  : GPIO_PIN_SET;  /* Pixel OFF */

    HAL_GPIO_WritePin(GPIOB, row_pin, row_state);
  }

  /* Select the column using the definitions in main.h. */
  switch (index) {
  case 0:
    selected_column = ENM0_Pin; /* PA2 */
    break;

  case 1:
    selected_column = ENM1_Pin; /* PA3 */
    break;

  case 2:
    selected_column = ENM2_Pin; /* PA10 */
    break;

  case 3:
    selected_column = ENM3_Pin; /* PA11 */
    break;

  case 4:
    selected_column = ENM4_Pin; /* PA12 */
    break;

  case 5:
    selected_column = ENM5_Pin; /* PA13 */
    break;

  case 6:
    selected_column = ENM6_Pin; /* PA14 */
    break;

  case 7:
    selected_column = ENM7_Pin; /* PA15 */
    break;

  default:
    return;
  }

  /* Enable only the selected column: active LOW. */
  HAL_GPIO_WritePin(GPIOA, selected_column, GPIO_PIN_RESET);
}
\end{stmcode}

Both the all-column mask and the selection
cases use \texttt{ENM0\_Pin}--\texttt{ENM7\_Pin}.
Column 0 selects PA2, and column 1 selects PA3.

The row loop uses consecutive GPIO pins
PB8--PB15. Disabling every column before
changing these outputs prevents the previous
column from briefly displaying the next
column's pattern.

\subsubsection{Complete application initialization}

After the existing GPIO and TIM2 initialization,
the application prepares the clock, initializes
the three software timers, and starts TIM2.
It then selects the initial matrix column
and records the matrix-scan timestamp.

\begin{stmcode}{Exercise 9 initialization inside main}
/* USER CODE BEGIN 2 */
HAL_GPIO_WritePin(GPIOA, GPIO_PIN_4 | GPIO_PIN_5, GPIO_PIN_SET);

updateClockBuffer();
update7SEG(index_led); /* Show the first digit immediately */

setTimer0(1000); /* PA5's first toggle */
setTimer1(1000); /* Next clock second and DOT toggle */
setTimer2(250);  /* Next display selection */

HAL_TIM_Base_Start_IT(&htim2);

updateLEDMatrix(index_led_matrix);
last_matrix_scan = HAL_GetTick();
/* USER CODE END 2 */
\end{stmcode}

The initial matrix index is zero, so the
first selected column is ENM0 on PA2.

\subsubsection{Complete main loop}

The complete loop below includes all inherited
tasks and the new matrix-scanning task.
The red-LED, clock, DOT, and seven-segment
processing follows Exercise 8.

\begin{stmcode}{Exercise 9 complete main loop}
/* USER CODE BEGIN WHILE */
while (1) {
    /* Red LED: first toggle after 1 s, then every 2 s */
    if (timer0_flag == 1) {
      HAL_GPIO_TogglePin(GPIOA, GPIO_PIN_5);
      setTimer0(2000);
    }

    /* Digital clock and DOT: update every 1 s */
    if (timer1_flag == 1) {
      setTimer1(1000);

      second++;

      if (second >= 60) {
        second = 0;
        minute++;
      }

      if (minute >= 60) {
        minute = 0;
        hour++;
      }

      if (hour >= 24) {
        hour = 0;
      }

      updateClockBuffer();
      HAL_GPIO_TogglePin(GPIOA, GPIO_PIN_4);
    }

    /* Seven-segment scan: change display every 250 ms */
    if (timer2_flag == 1) {
      setTimer2(250);

      index_led++;
      if (index_led >= MAX_LED) {
        index_led = 0;
      }

      update7SEG(index_led);
    }

    /* Exercise 09: change matrix column every 5 ms */
    uint32_t now = HAL_GetTick();

    if ((uint32_t)(now - last_matrix_scan) >= 5U) {
      last_matrix_scan = now;

      index_led_matrix++;
      if (index_led_matrix >= MAX_LED_MATRIX) {
        index_led_matrix = 0;
      }

      updateLEDMatrix(index_led_matrix);
    }
    /* USER CODE END WHILE */

    /* USER CODE BEGIN 3 */
}
/* USER CODE END 3 */
\end{stmcode}

Timer 0 controls PA5, timer 1 controls the
clock and DOT, and timer 2 controls
seven-segment display selection.

Matrix scanning uses the HAL millisecond
time base. The comparison with \texttt{5U}
selects a nominal 5-ms column interval.
The index advances through columns 0--7
and returns to zero after column 7.

The stored A pattern remains unchanged.
The GPIO outputs retain each selected
column and its row pattern between updates.
No \texttt{HAL\_Delay()} is used.

\subsubsection{Expected behavior}

The smaller A remains stationary while
its columns are repeatedly scanned.

The clock advances approximately once
per second, DOT changes state every second,
and the selected seven-segment display
changes every 250~ms.

PA5 starts off, toggles after approximately
one second, and then changes state every
two seconds.


% ============================================================
\subsection{Exercise 10: Scroll the matrix character to the left}

\subsubsection{Objective and inherited configuration}

Exercise 10 adds left-scrolling movement
to the A displayed in Exercise 9.

The circuit remains as shown in
Figure~\ref{fig:ex09}. GPIO assignments
match Table~\ref{tab:ex09-matrix-pins}:
ENM0 uses PA2, ENM1 uses PA3, and
ENM2--ENM7 use PA10--PA15.

The \texttt{updateLEDMatrix()} function,
initial matrix buffer, clock functions,
software-timer functions, GPIO initialization,
and \texttt{HAL\_TIM\_PeriodElapsedCallback()}
match exactly the corresponding parts
in Exercise 9.

Column scanning retains its 5-ms interval.
The new animation uses a nominal 200-ms
movement interval.

\subsubsection{Additional animation variables}

The matrix declarations from Exercise 9
are retained. The following declarations
are added inside \texttt{USER CODE BEGIN PV}.

\begin{stmcode}{Exercise 10 animation variables and source pattern}
uint32_t last_matrix_animation = 0;
int matrix_scroll_offset = 0;

/* Small A with margins, followed by eight blank columns */
const uint8_t scrolling_pattern[16] = {
    0x00, 0x7C, 0x12, 0x12,
    0x12, 0x7C, 0x00, 0x00,

    0x00, 0x00, 0x00, 0x00,
    0x00, 0x00, 0x00, 0x00
};
\end{stmcode}

The first eight source bytes match the
initial A frame from Exercise 9.
The next eight bytes provide blank columns
between successive appearances of the character.

The matrix displays an eight-column window
from this 16-column source.
\texttt{matrix\_scroll\_offset} determines
the starting position of that window.

\subsubsection{Report 1: Left-scrolling function}

The new function \texttt{scrollMatrixLeft()}
is defined inside \texttt{USER CODE BEGIN 0}.

\begin{stmcode}{Exercise 10 scrollMatrixLeft function}
void scrollMatrixLeft(void) {
  matrix_scroll_offset++;

  if (matrix_scroll_offset >= 16) {
    matrix_scroll_offset = 0;
  }

  for (int col = 0; col < MAX_LED_MATRIX; col++) {
    int source = (col + matrix_scroll_offset) % 16;
    matrix_buffer[col] = scrolling_pattern[source];
  }
}
\end{stmcode}

Each call advances the viewing-window offset
by one and rebuilds all eight display columns.
Each physical column receives the next source
column, making the visible character move left.

The modulo operation wraps the viewing
window around the 16-column source.
After 16 movements, the initial position repeats.

The function changes the buffer contents.
The inherited \texttt{updateLEDMatrix()}
function displays those contents using
the updated ENM pin definitions.

\subsubsection{Complete application initialization}

Initialization follows Exercise 9, with
the animation timestamp initialized from
the matrix-scan timestamp.

\begin{stmcode}{Exercise 10 initialization inside main}
/* USER CODE BEGIN 2 */
HAL_GPIO_WritePin(GPIOA, GPIO_PIN_4 | GPIO_PIN_5, GPIO_PIN_SET);

updateClockBuffer();
update7SEG(index_led); /* Show the first digit immediately */

setTimer0(1000); /* PA5's first toggle */
setTimer1(1000); /* Next clock second and DOT toggle */
setTimer2(250);  /* Next display selection */

HAL_TIM_Base_Start_IT(&htim2);

updateLEDMatrix(index_led_matrix);
last_matrix_scan = HAL_GetTick();

last_matrix_animation = last_matrix_scan;
/* USER CODE END 2 */
\end{stmcode}

The scroll offset initially equals zero.
The initial buffer therefore shows the same
A position as Exercise 9.

\subsubsection{Complete main loop}

The complete loop retains all inherited
tasks and adds animation inside the
matrix-scanning block.

\begin{stmcode}{Exercise 10 complete main loop}
/* USER CODE BEGIN WHILE */
while (1) {
    /* Red LED: first toggle after 1 s, then every 2 s */
    if (timer0_flag == 1) {
      HAL_GPIO_TogglePin(GPIOA, GPIO_PIN_5);
      setTimer0(2000);
    }

    /* Digital clock and DOT: update every 1 s */
    if (timer1_flag == 1) {
      setTimer1(1000);

      second++;

      if (second >= 60) {
        second = 0;
        minute++;
      }

      if (minute >= 60) {
        minute = 0;
        hour++;
      }

      if (hour >= 24) {
        hour = 0;
      }

      updateClockBuffer();
      HAL_GPIO_TogglePin(GPIOA, GPIO_PIN_4);
    }

    /* Seven-segment scan: change display every 250 ms */
    if (timer2_flag == 1) {
      setTimer2(250);

      index_led++;
      if (index_led >= MAX_LED) {
        index_led = 0;
      }

      update7SEG(index_led);
    }

    uint32_t now = HAL_GetTick();

    /* Scan one matrix column every 5 ms */
    if ((uint32_t)(now - last_matrix_scan) >= 5U) {
      last_matrix_scan = now;

      index_led_matrix++;
      if (index_led_matrix >= MAX_LED_MATRIX) {
        index_led_matrix = 0;
      }

      /* Move A at the beginning of a complete scan */
      if (index_led_matrix == 0 &&
          (uint32_t)(now - last_matrix_animation) >= 200U) {

        last_matrix_animation += 200U;
        scrollMatrixLeft();
      }

      updateLEDMatrix(index_led_matrix);
    }
    /* USER CODE END WHILE */

    /* USER CODE BEGIN 3 */
}
/* USER CODE END 3 */
\end{stmcode}

The red-LED, clock, DOT, and seven-segment
blocks retain the behavior of Exercise 9.

The outer matrix condition controls the
5-ms column interval. The inner condition
allows movement when the animation deadline
has been reached and the next displayed
column is column 0.

The buffer is updated immediately before
displaying column 0 through ENM0 on PA2.
It remains unchanged while the remaining
columns of that scan are displayed.
This prevents the character position from
changing halfway through a complete scan.

The statement
\texttt{last\_matrix\_animation += 200U}
advances the scheduled movement time
by a fixed interval. Actual movement occurs
at the first eligible beginning of a scan.

Scanning and animation both remain in
the main loop. The TIM2 callback continues
calling only \texttt{timer\_run()}.

\subsubsection{Expected behavior}

The smaller A begins in the same position
as in Exercise 9. It moves left one column
at a time, leaves the display, passes through
the blank portion of the source, and re-enters
from the right.

The character keeps its original shape.
The clock, DOT, PA5 red LED, and seven-segment
tasks continue operating during the animation,
without a blocking delay.

\clearpage
\phantomsection
\addcontentsline{toc}{section}{References}
\begin{thebibliography}{9}
\raggedright

\bibitem{lab02-repository}
V. L. T. Trieu,
\textit{Lab 02 code + schematic}, GitHub repository. [Online]. Available:
\url{https://github.com/volethanhtrieu/Lab-exercise-MPU---MCU/tree/main/lab02}.

\bibitem{lab02-manual}
Ho Chi Minh City University of Technology,
\textit{Microcontroller Laboratory -- Lab 02: Timer Interrupt and LED Scanning},
course laboratory manual.

\bibitem{stm32f103-datasheet}
STMicroelectronics,
\textit{STM32F103x4 and STM32F103x6 Low-density Performance Line Datasheet},
DocID15060, Rev. 7, June 2015. [Online]. Available:
\url{https://www.st.com/resource/en/datasheet/stm32f103c4.pdf}.
[Accessed: September 18, 2026].

\bibitem{stm32cubemx}
STMicroelectronics,
\textit{STM32CubeMX: STM32Cube Initialization Code Generator}. [Online].
Available: \url{https://www.st.com/en/development-tools/stm32cubemx.html}.
[Accessed: September 18, 2026].

\bibitem{stm32cubef1}
STMicroelectronics,
\textit{STM32CubeF1 Embedded Software Package}. [Online]. Available:
\url{https://www.st.com/en/embedded-software/stm32cubef1.html}.
[Accessed: September 18, 2026].

\bibitem{proteus-simulation}
Labcenter Electronics,
\textit{Proteus Circuit Simulation Software with SPICE}. [Online]. Available:
\url{https://www.labcenter.com/simulation/}.
[Accessed: September 18, 2026].

\end{thebibliography}
\end{document}
